\documentclass[11pt]{article}

\usepackage{amsmath}
\usepackage{amsfonts}
\usepackage{amsthm,thmtools}
\usepackage{tikz}
\usepackage{enumerate}


\newcommand{\batsym}{\begin{tikzpicture}[scale=0.05]
	\draw[fill,black]   (-0.25,1.48) .. controls (-0.1,1.5) and (0.1,1.5) .. (0.25,1.48) -- (0.35,1.92) .. controls (0.425,1.8) and (0.41,1.3) .. (0.45,1.2) .. controls (0.6,1.05) and (1.96,1.05) .. (1.98,2.08) -- (5.93,2.08) .. controls (4.2,1.45) and (4,0.3) .. (4.2,-0.28) .. controls (2.4,-0.09) and (0.4,-0.5) .. (0,-2.052) .. controls (-0.4,-0.5) and (-2.4,-0.09) .. (-4.2,-0.28) .. controls (-4,0.3) and (-4.2,1.45) .. (-5.93,2.08) -- (-1.98,2.08) .. controls (-1.96,1.05) and (-0.6,1.05) .. (-0.45,1.2) .. controls (-0.41,1.3) and (-0.425,1.8) .. (-0.35,1.92) -- (-0.25,1.48);
	\end{tikzpicture}}
\newcommand{\QEB}{\begin{flushright}\batsym\end{flushright}}



\def\bold{\textbf}
\theoremstyle{plain}
\newtheorem{problem}{Problem}
\theoremstyle{remark}
\newtheorem*{idea}{Idea}

\renewcommand*{\proofname}{Solution}


\renewcommand{\qedsymbol}{$\batsym$}

\newcommand {\ds}{\displaystyle\sum}
\newcommand {\nnR} {\mathbb{R}_{\geq 0}}

\newcommand {\Mtt}[4] {\begin{bmatrix}#1 & #2\\ #3 & #4\end{bmatrix}}
\newcommand {\Ctt}[2] {\Mtt {#1}{-#2}{#2}{#1}}
\newcommand {\mbR} {\mathbb R}
\newcommand {\mcC} {\mathcal C}
\newcommand {\mbZ} {\mathbb Z}
\newcommand {\mbN} {\mathbb N}
\newcommand {\mbQ} {\mathbb Q}
\newcommand {\RnQ} {\mathbb R \setminus \mathbb Q}

\begin{document}
	\title{Math 3A03 - Tutorial 2 Questions - Winter 2019}
	\author{Nikolay Hristov}
	\date{January 21/23, 2019}
	\maketitle
	
	\begin{problem}
		Use the induction principle to prove that $\displaystyle \sum_{k=1}^n k^2 = \frac{n(n+1)(2n+1)}{6}$ for any $n \in \mbN$.
	\end{problem}

	\begin{problem}
		\label{density_of_set}
		Show that the set $P:=\left\{\sqrt{2}q : q \in \mathbb{Q}\right\}$ is dense in $\mathbb{R}$.
	\end{problem}
	
	\begin{problem}
		If a set $E$ is dense in $\mathbb{R}$, what can you conclude about
		\begin{enumerate}[(a)]
			\item a set $A$, such that $E \subset A$?
			
			\item the set $\mathbb{R} \backslash E$?
			
			\item the set $E \cap F$, where $F$ is also dense in $\mathbb{R}$?
		\end{enumerate}
	\end{problem}

	\begin{problem}
	Let $I := (a,b) \subseteq \mathbb{R}$ be an interval with $b \in \mathbb{R}$, $a<b$, and $E \subseteq I$ be a dense subset of I.
	
	\begin{enumerate}[(a)]
		\item Show that sup$(E)=$sup$(I) = b$. 
		\item Prove that $\sqrt{5}$ is the sup$\left( [0, \sqrt{5}) \right)$.
		\item Prove that $\sqrt{5}$ is 
		the sup$\left( [0, \sqrt{5}] \cap	\mathbb{Q} \right)$.
	\end{enumerate}
	\end{problem}

	\begin{problem}
		Prove, using the definition of the limit $\displaystyle \lim_{n\rightarrow \infty} \frac {\cos(n)}{\sqrt{n^5+n^2+3 }}= 0$.
	\end{problem}
	
\end{document}